\section{Derivations and interpretation}
\subsection{What SNAP estimates}
SNAP abbreviates Seed Noise Across Phenotypes, with a phenotype denoting a benchmark score. Its target is the covariance of run deviations that generalise across item halves. DataDecide replicates vary initialisation and data order together, and auxiliary runs below 1B use shortened training schedules. Common-step scoring aligns the observation step, while schedule differences remain within the full-set target. We therefore use run covariance for the empirical estimate, reserving a seed-only interpretation for designs with matching training conditions.

The distinction between item-general variation and fixed-battery variation determines how a practitioner can use the estimate. A run can perform differently on particular items without changing its expected score over an item population. Under the measurement model, that variation enters the item error and doesn't contribute to the cross-half covariance target. Repeated evaluation of a fixed item bank can retain this component, so its run variance needn't equal the item-general variance estimated here.

The analogy with quantitative genetics distinguishes covariance between configurations from residual covariance within them. It doesn't identify separate additive-genetic and shared-environment components in a model-training experiment. Between-configuration correlation $R_P$ reflects differences in recipe and model size, while the within-configuration covariance $\Sig$ concerns replicate deviations. The plug-in comparison asks whether the first correlation structure predicts the second after supplying marginal run variances.

\subsection{Notation}
\begin{table}[h]
\centering
\caption{Notation for the cross-half construction and its empirical diagnostics. The coefficient $\bar r_E$ uses covariance normalisation, which differs from mean correlation under unequal marginal variances.}
\begin{tabular}{lp{0.72\linewidth}}
\toprule
Symbol & Meaning\\
\midrule
$c,r,j,i$ & Configuration, replicate run, benchmark, and item indices\\
$N,R,K$ & Configuration, replicate, and benchmark counts, with values 125, 3, and 10\\
$\Sigma_{E,c},\Sig$ & Configuration covariance and its average across the design\\
$R_E,R_P$ & Run-effect correlation and between-configuration score correlation\\
$y^M,y^A$ & Mean per-byte margin and accuracy on the specified item set\\
$d^H_{crj},g^H_{cr}$ & Centred benchmark score and its benchmark average on half $H$\\
$T_c,U_c$ & Cross-half aggregate covariance and its diagonal-only contribution\\
$\Lam,\bar r_E$ & Standard-deviation inflation and effective covariance coefficient\\
$K_{\mathrm{eff}}$ & Independent-benchmark equivalent at the same marginal RMS deviation\\
$\sigma_{\mathrm{agg}},\sigma_{\mathrm{indep}}$ & Aggregate deviation with full covariance and under independence\\
$\rho_g,\bar v$ & Aggregate reliability and its scaled item-noise contribution\\
$s_{cr},x_{cr}^{(-j)}$ & Gain covariate and opposite-half competence proxy\\
\bottomrule
\end{tabular}
\end{table}

\subsection{Cross-half identity and its assumptions}
Fix a configuration and suppress its index for clarity. Write $y^H_{rj}=\mu^H_j+E_{rj}+\varepsilon^H_{rj}$, so centring across runs gives
\begin{equation}
d^H_{rj}=(E_{rj}-\bar E_{\cdot j})+(\varepsilon^H_{rj}-\bar\varepsilon^H_{\cdot j}).
\end{equation}
Assume that run-effect vectors are independent draws with covariance $\Sig$. Conditional on those effects, assume zero-mean item errors and zero covariance between $\varepsilon^A_{rj}$ and $\varepsilon^B_{sk}$ for the run and benchmark pairs in the calculation. Expanding $d^A_{rj}d^B_{rk}$ then removes its error cross terms in expectation, leaving
\begin{align}
\mathbb E[d^A_{rj}d^B_{rk}]
&=\mathbb E[(E_{rj}-\bar E_{\cdot j})(E_{rk}-\bar E_{\cdot k})]\\
&=(1-1/R)\Sigma_E(j,k).
\end{align}
Summing across runs and dividing by $R-1$ gives the unbiased covariance moment. We symmetrise the two half orderings for the matrix estimate,
\begin{equation}
\widehat\Sigma_E(j,k)=\frac{1}{2N(R-1)}\sum_{c,r}
\left(d^A_{crj}d^B_{crk}+d^A_{crk}d^B_{crj}\right).
\end{equation}
Symmetry follows from this construction, while positive semidefiniteness doesn't follow for a finite sample. A negative diagonal moment near zero therefore doesn't by itself falsify the identity, although it prevents ordinary correlation normalisation for that trait.

The independence assumption concerns item effects, rather than whether two identifiers differ. Shared source passages, duplicate questions, and overlapping content can create cross-half covariance. A fixed finite item bank split without replacement also requires a model for its sampling dependence. The available item-half assignment alone can't verify these conditions.

BoolQ is a specific unresolved case in this analysis. Questions associated with one passage can enter both halves, and a run-level response to that passage can correlate their errors. The original analysis didn't join public passage text to the retained item identifiers. Its covariance estimate can therefore include a passage-dependent diagonal contribution that the item-general target excludes. Removing BoolQ changes both this exposure and the benchmark composition, so the removal sweep can't identify the size of the leakage alone.

For a diagonal additional covariance $D$ with nonnegative trace, the inflation ratio becomes
\begin{equation}
\Lam_D^2=\frac{\mathbf1^{\mathsf T}\Sig\mathbf1+\operatorname{tr}D}
{\operatorname{tr}\Sig+\operatorname{tr}D}.
\end{equation}
This ratio lies between $\Lam^2$ and one when the denominators are positive. It explains attenuation toward one under that restricted model, but doesn't yield a general upper or lower bound when shared-item effects also contribute off-diagonal covariance. The same restriction is needed before using the item-general factor to bound fixed-battery inflation.

\subsection{Covariance coefficient and effective count}
Let $\sigma_j^2=\Sigma_E(j,j)$ and $r_{jk}=\Sigma_E(j,k)/(\sigma_j\sigma_k)$ whenever both variances are positive. The finite-battery identity is
\begin{equation}
\Lam^2=1+\frac{\sum_{j\ne k}r_{jk}\sigma_j\sigma_k}{\sum_j\sigma_j^2}
=1+(K-1)\bar r_E.
\end{equation}
The normalised weighted mean correlation instead divides the numerator by $\sum_{j\ne k}\sigma_j\sigma_k$. If that denominator is positive, then
\begin{equation}
\bar r_E=a_\sigma\bar r_{\mathrm{weighted}},\qquad
a_\sigma=\frac{\sum_{j\ne k}\sigma_j\sigma_k}{(K-1)\sum_j\sigma_j^2}.
\end{equation}
For nonnegative standard deviations, $a_\sigma$ lies between zero and one. With two perfectly correlated traits having standard deviations one and two, the effective coefficient is 0.8 even though the weighted mean correlation is one. Equal marginal variances make the coefficient and the mean correlation coincide.

Writing $\bar\sigma_E^2=K^{-1}\operatorname{tr}\Sig$, the aggregate variance satisfies
\begin{equation}
\sigma_{\mathrm{agg}}^2=\bar\sigma_E^2\left[\frac1K+\frac{K-1}{K}\bar r_E\right].
\end{equation}
A nonzero asymptotic floor follows only for a sequence of batteries whose average variance and effective coefficient converge to suitable limits. Bounded variances alone don't imply that convergence, and the present ten-benchmark study doesn't estimate such a sequence. We use the finite-battery identity for the reported effective counts.

For positive aggregate variance, $K_{\mathrm{eff}}=K/\Lam^2$ compares the measured average with an independent average having the same marginal RMS deviation. Negative aggregate off-diagonal covariance can give $K_{\mathrm{eff}}>K$. A single positive-variance benchmark has $\Lam=1$, although its off-diagonal coefficient is undefined because $K-1=0$. This arithmetic count also differs from a spectral participation ratio.

\subsection{Pooling and finite-sample bias}
We pool $T_c$ and $U_c$ across configurations before taking their ratio. Each individual denominator uses two replicate degrees of freedom and can approach zero or become negative. Averaging individual ratios would give those unstable denominators disproportionate influence. Pooling reduces that problem but doesn't remove the effect of configurations with large covariance contributions.

The reported recipe influence participation ratios are 2.88 for margins and 8.68 for accuracy. One recipe contributes 0.543 of the squared margin influence, while removing that recipe moves the estimate by at most 0.037. These quantities describe influence concentration, and aren't literal counts of independent clusters for inference. The reported permutation coverage of 0.925 on margins is a reason to qualify the nominal interval, rather than evidence that a bootstrap automatically resolves concentration.

The covariance moments are unbiased under their assumptions, while the pooled ratio has a random denominator. Its square root introduces another nonlinear transformation, so neither transformation inherits exact unbiasedness. A finite implementation must also specify what happens when a pooled sum is nonpositive. The reported data give finite full-battery estimates, but that observation doesn't establish finite estimates for arbitrary input populations.

\subsection{Permutation reference}
For each configuration and trait, apply a uniform permutation $\pi_j$ of run labels to both halves together. The diagonal products then sum to the same value, leaving $U_c$ unchanged. Independent permutations across traits give zero expected off-diagonal products because each centred deviation sums to zero. Conditional on a positive fixed pooled denominator, this implies
\begin{equation}
\mathbb E_\pi\left[\frac{\sum_cT_c^\pi}{\sum_cU_c}\right]=1.
\end{equation}
It doesn't imply that $\mathbb E_\pi[\widehat\Lam^\pi]=1$. If all permuted ratios are nonnegative, concavity of the square root gives an expectation no larger than one, with strict inequality for a nonconstant ratio. For example, equally likely ratios of 0.5 and 1.5 average to one, while their square roots average to about 0.966.

A common permutation across all traits preserves their pairings and can't remove cross-trait association. Independent trait permutations avoid that problem, and the identity belongs in the uniform permutation sample space. Its chance of reproducing the observation doesn't invalidate a permutation procedure. The construction preserves marginal score values and cluster membership, but can alter dependence of run alignments across configurations, so validity as a test requires an exchangeability assumption beyond zero mean covariance.

\subsection{Planning calculations and information ratio}
The author-reported power calculation assumes approximately Gaussian half-score deviations. With $d^A=e+a$, $d^B=e+b$, independent errors $a,b$, and $\operatorname{Var}(a)=\operatorname{Var}(b)=k\operatorname{Var}(e)$, the product variance is
\begin{equation}
\operatorname{Var}(d^Ad^B)=\sigma^4[(1+k)^2+1],\qquad
k=2(1-\rho_g)/\rho_g.
\end{equation}
The planning values used two replicate degrees of freedom per configuration, numerator-denominator correlation 0.7, and design effect 1.60. Appendix B preserves the resulting values for the intended 85-configuration analysis, which differs from the realised 125-configuration study. We don't infer realised power from the observed effect.

For the information calculation, let item margins follow a Gaussian location model with standard deviation $s_m$. An infinitesimal shift $\delta$ changes accuracy by $f(0)\delta+o(\delta)$, while it changes the mean margin by $\delta$. Under independent item sampling, the item-noise variances are $p(1-p)/n$ and $s_m^2/n$, and the Gaussian model gives $f(0)=\phi(z_0)/s_m$. Dividing noise by squared local sensitivity yields Equation~\ref{eq:information}. The scale cancellation and minimum $\pi/2$ belong to that model, rather than arbitrary empirical margin distributions.

The competence and gain adjustments have additional measurement limitations. The opposite-half competence proxy uses estimated seed standard deviations, including a $10^{-6}$ floor where the margin variance estimate for WinoGrande is $-3.15\times10^{-7}$. The gain covariate uses both halves, so its item noise can correlate with the outcome errors. The reported zero-mediation simulation treats this shared item component differently from the observed regression. Its margin reference moves from 1.24502 to 1.0910, with standard deviation 0.096, while the fully adjusted observed value is 1.276. The accuracy reference begins at 1.11831 rather than the observed 1.07837, and ends at 1.1045 against an adjusted 1.094. These mismatches prevent a calibrated causal interpretation of the residualisation comparisons.
